% Options for packages loaded elsewhere
\PassOptionsToPackage{unicode}{hyperref}
\PassOptionsToPackage{hyphens}{url}
\documentclass[
  10pt,
]{article}
\usepackage{xcolor}
\usepackage[margin=1in]{geometry}
\usepackage{amsmath,amssymb}
\setcounter{secnumdepth}{-\maxdimen} % remove section numbering
\usepackage{iftex}
\ifPDFTeX
  \usepackage[T1]{fontenc}
  \usepackage[utf8]{inputenc}
  \usepackage{textcomp} % provide euro and other symbols
\else % if luatex or xetex
  \usepackage{unicode-math} % this also loads fontspec
  \defaultfontfeatures{Scale=MatchLowercase}
  \defaultfontfeatures[\rmfamily]{Ligatures=TeX,Scale=1}
\fi
\usepackage{lmodern}
\ifPDFTeX\else
  % xetex/luatex font selection
\fi
% Use upquote if available, for straight quotes in verbatim environments
\IfFileExists{upquote.sty}{\usepackage{upquote}}{}
\IfFileExists{microtype.sty}{% use microtype if available
  \usepackage[]{microtype}
  \UseMicrotypeSet[protrusion]{basicmath} % disable protrusion for tt fonts
}{}
\makeatletter
\@ifundefined{KOMAClassName}{% if non-KOMA class
  \IfFileExists{parskip.sty}{%
    \usepackage{parskip}
  }{% else
    \setlength{\parindent}{0pt}
    \setlength{\parskip}{6pt plus 2pt minus 1pt}}
}{% if KOMA class
  \KOMAoptions{parskip=half}}
\makeatother
\usepackage{color}
\usepackage{fancyvrb}
\newcommand{\VerbBar}{|}
\newcommand{\VERB}{\Verb[commandchars=\\\{\}]}
\DefineVerbatimEnvironment{Highlighting}{Verbatim}{commandchars=\\\{\}}
% Add ',fontsize=\small' for more characters per line
\newenvironment{Shaded}{}{}
\newcommand{\AlertTok}[1]{\textcolor[rgb]{1.00,0.00,0.00}{\textbf{#1}}}
\newcommand{\AnnotationTok}[1]{\textcolor[rgb]{0.38,0.63,0.69}{\textbf{\textit{#1}}}}
\newcommand{\AttributeTok}[1]{\textcolor[rgb]{0.49,0.56,0.16}{#1}}
\newcommand{\BaseNTok}[1]{\textcolor[rgb]{0.25,0.63,0.44}{#1}}
\newcommand{\BuiltInTok}[1]{\textcolor[rgb]{0.00,0.50,0.00}{#1}}
\newcommand{\CharTok}[1]{\textcolor[rgb]{0.25,0.44,0.63}{#1}}
\newcommand{\CommentTok}[1]{\textcolor[rgb]{0.38,0.63,0.69}{\textit{#1}}}
\newcommand{\CommentVarTok}[1]{\textcolor[rgb]{0.38,0.63,0.69}{\textbf{\textit{#1}}}}
\newcommand{\ConstantTok}[1]{\textcolor[rgb]{0.53,0.00,0.00}{#1}}
\newcommand{\ControlFlowTok}[1]{\textcolor[rgb]{0.00,0.44,0.13}{\textbf{#1}}}
\newcommand{\DataTypeTok}[1]{\textcolor[rgb]{0.56,0.13,0.00}{#1}}
\newcommand{\DecValTok}[1]{\textcolor[rgb]{0.25,0.63,0.44}{#1}}
\newcommand{\DocumentationTok}[1]{\textcolor[rgb]{0.73,0.13,0.13}{\textit{#1}}}
\newcommand{\ErrorTok}[1]{\textcolor[rgb]{1.00,0.00,0.00}{\textbf{#1}}}
\newcommand{\ExtensionTok}[1]{#1}
\newcommand{\FloatTok}[1]{\textcolor[rgb]{0.25,0.63,0.44}{#1}}
\newcommand{\FunctionTok}[1]{\textcolor[rgb]{0.02,0.16,0.49}{#1}}
\newcommand{\ImportTok}[1]{\textcolor[rgb]{0.00,0.50,0.00}{\textbf{#1}}}
\newcommand{\InformationTok}[1]{\textcolor[rgb]{0.38,0.63,0.69}{\textbf{\textit{#1}}}}
\newcommand{\KeywordTok}[1]{\textcolor[rgb]{0.00,0.44,0.13}{\textbf{#1}}}
\newcommand{\NormalTok}[1]{#1}
\newcommand{\OperatorTok}[1]{\textcolor[rgb]{0.40,0.40,0.40}{#1}}
\newcommand{\OtherTok}[1]{\textcolor[rgb]{0.00,0.44,0.13}{#1}}
\newcommand{\PreprocessorTok}[1]{\textcolor[rgb]{0.74,0.48,0.00}{#1}}
\newcommand{\RegionMarkerTok}[1]{#1}
\newcommand{\SpecialCharTok}[1]{\textcolor[rgb]{0.25,0.44,0.63}{#1}}
\newcommand{\SpecialStringTok}[1]{\textcolor[rgb]{0.73,0.40,0.53}{#1}}
\newcommand{\StringTok}[1]{\textcolor[rgb]{0.25,0.44,0.63}{#1}}
\newcommand{\VariableTok}[1]{\textcolor[rgb]{0.10,0.09,0.49}{#1}}
\newcommand{\VerbatimStringTok}[1]{\textcolor[rgb]{0.25,0.44,0.63}{#1}}
\newcommand{\WarningTok}[1]{\textcolor[rgb]{0.38,0.63,0.69}{\textbf{\textit{#1}}}}
\usepackage{longtable,booktabs,array}
\usepackage{calc} % for calculating minipage widths
% Correct order of tables after \paragraph or \subparagraph
\usepackage{etoolbox}
\makeatletter
\patchcmd\longtable{\par}{\if@noskipsec\mbox{}\fi\par}{}{}
\makeatother
% Allow footnotes in longtable head/foot
\IfFileExists{footnotehyper.sty}{\usepackage{footnotehyper}}{\usepackage{footnote}}
\makesavenoteenv{longtable}
\setlength{\emergencystretch}{3em} % prevent overfull lines
\providecommand{\tightlist}{%
  \setlength{\itemsep}{0pt}\setlength{\parskip}{0pt}}
\usepackage{bookmark}
\IfFileExists{xurl.sty}{\usepackage{xurl}}{} % add URL line breaks if available
\urlstyle{same}
\hypersetup{
  hidelinks,
  pdfcreator={LaTeX via pandoc}}

\author{}
\date{}

\begin{document}

\section{PlatePlanner: An AI-Powered Meal Planning System with
Graph-Based Ingredient Substitution, Food Ontology Reasoning, and
Ethnicity-Aware Recipe
Recommendation}\label{plateplanner-an-ai-powered-meal-planning-system-with-graph-based-ingredient-substitution-food-ontology-reasoning-and-ethnicity-aware-recipe-recommendation}

\textbf{IEEE CMPE 295B --- Master's Project Report}

\textbf{Authors:} - Sandilya Chimalamarri --- ML Research Lead, San José
State University - Sai Priyanka Bonkuri --- Graph Database Architect,
San José State University - Pavan Charith Devarapalli --- Backend
Systems Engineer, San José State University - Sai Dheeraj Gollu --- Data
Pipeline Engineer, San José State University

\textbf{Advisor:} {[}Project Advisor Name{]}, San José State University

\begin{center}\rule{0.5\linewidth}{0.5pt}\end{center}

\subsection{Abstract}\label{abstract}

PlatePlanner is an AI-powered meal planning and recipe recommendation
platform that integrates multiple machine learning models to deliver
personalized, culturally-aware, and nutritionally sound meal plans. The
system combines (1) a FAISS-based semantic retrieval engine using
fine-tuned sentence transformers, (2) a Graph Neural Network (GraphSAGE)
trained on a Neo4j ingredient knowledge graph for context-aware
ingredient substitution, (3) a food ontology layer encoding culinary
rules and dietary constraints, and (4) a Retrieval-Augmented Generation
(RAG) pipeline that leverages large language models for recipe
adaptation. Evaluated on the RecipeNLG dataset (2.2M+ recipes) augmented
with international cuisine corpora, PlatePlanner achieves a GNN
substitution AUC of 0.87, semantic retrieval latency under 100ms at p95,
92\% user satisfaction on context-aware substitutions, and dietary
classification accuracy exceeding 90\% across 12 diet types and 11
allergen categories.

\textbf{Keywords:} Recipe Recommendation, Graph Neural Networks,
Knowledge Graphs, Ingredient Substitution, Retrieval-Augmented
Generation, Dietary Classification, Food Ontology, Meal Planning.

\begin{center}\rule{0.5\linewidth}{0.5pt}\end{center}

\subsection{I. Introduction}\label{i.-introduction}

The intersection of artificial intelligence and personal nutrition
represents one of the most impactful application domains for modern ML
systems. Despite the proliferation of recipe apps, most existing
platforms rely on keyword search and manual filtering, failing to
leverage the rich structural relationships between ingredients,
cuisines, and dietary needs {[}1{]}.

PlatePlanner addresses four key research challenges:

\begin{enumerate}
\def\labelenumi{\arabic{enumi}.}
\tightlist
\item
  \textbf{Semantic Recipe Retrieval} --- Finding relevant recipes given
  a set of pantry ingredients, without requiring exact keyword matches.
\item
  \textbf{Context-Aware Ingredient Substitution} --- Recommending
  culinary-valid replacements for missing ingredients using
  graph-structural reasoning.
\item
  \textbf{Food Ontology and Safety Filtering} --- Deterministically
  enforcing dietary restrictions and allergen constraints through a
  Neo4j-backed ontology.
\item
  \textbf{Ethnicity and Preference-Based Filtering} --- Supporting
  cuisine preferences (Mexican, Indian, Italian, Asian, Mediterranean,
  American, and more) combined with dietary profiles.
\end{enumerate}

The system is deployed as a FastAPI microservice consumed by a React
Native mobile application, with Railway cloud hosting and Neon
PostgreSQL as the production database.

\begin{center}\rule{0.5\linewidth}{0.5pt}\end{center}

\subsection{II. Related Work}\label{ii.-related-work}

\textbf{Recipe Retrieval.} Early systems used TF-IDF and BM25 over
recipe corpora {[}2{]}. More recently, dense retrieval with bi-encoder
models such as DPR {[}3{]} and SentenceTransformers {[}4{]} has become
standard. PlatePlanner extends these with a task-specific fine-tuned
encoder on recipe ingredient sequences.

\textbf{Ingredient Substitution.} Rule-based approaches {[}5{]} rely on
manually curated substitution tables. Stanton et al.~{[}6{]} used
ingredient embeddings from co-occurrence statistics. More recent work
uses GNNs over knowledge graphs {[}7{]}. PlatePlanner combines Word2Vec
co-occurrence embeddings with GraphSAGE link prediction.

\textbf{Food Knowledge Graphs.} FoodKG {[}8{]} and FlavorGraph {[}9{]}
demonstrate the value of graph-structured food data. PlatePlanner
constructs a domain-specific Neo4j property graph with typed
relationships tailored for substitution and dietary reasoning.

\textbf{Dietary Classification.} Prior work uses ingredient lookup
tables {[}10{]} and NLP-based classification {[}11{]}. PlatePlanner
employs compound-ingredient-aware regex matching with exception lists to
avoid false positives (e.g., ``coconut milk'' should not trigger dairy
restrictions).

\textbf{Retrieval-Augmented Generation.} Lewis et al.~{[}12{]}
introduced RAG for open-domain QA. PlatePlanner applies RAG to recipe
adaptation, combining FAISS retrieval with LLM generation for
personalized cooking guidance.

\begin{center}\rule{0.5\linewidth}{0.5pt}\end{center}

\subsection{III. System Architecture}\label{iii.-system-architecture}

The PlatePlanner system is organized into four layers as shown in Figure
1.

\textbf{Fig. 1 --- PlatePlanner System Architecture} \emph{(See
generated system\_architecture diagram)}

\subsubsection{A. User Interface Layer}\label{a.-user-interface-layer}

A React Native / Expo mobile application communicates with a FastAPI
REST backend. The app supports authentication (JWT + Google OAuth),
pantry management, meal planning, shopping list generation, nutrition
tracking, and recipe browsing.

\subsubsection{B. Intelligence Layer}\label{b.-intelligence-layer}

Four ML/AI components form the core intelligence: - \textbf{FAISS
Semantic Search} --- Dense retrieval over 2.2M+ recipe embeddings -
\textbf{GNN Substitution Engine} --- GraphSAGE-based link prediction on
the ingredient graph - \textbf{Dietary Classifier} --- Rule-based NLP
for dietary and allergen classification - \textbf{RAG Service} ---
LLM-powered recipe adaptation using Gemini / GPT-4 / Llama3

\subsubsection{C. Data Layer}\label{c.-data-layer}

\begin{itemize}
\tightlist
\item
  \textbf{Neo4j Graph Database} --- Ingredient and recipe nodes with
  typed relationships
\item
  \textbf{PostgreSQL (Neon)} --- User accounts, meal plans, nutrition
  cache, USDA data
\item
  \textbf{SQLite} --- Local recipe index (recipes.db, 2.2M+ recipes)
\end{itemize}

\subsubsection{D. External APIs}\label{d.-external-apis}

\begin{itemize}
\tightlist
\item
  \textbf{USDA FoodData Central} --- Ground-truth nutrition data per
  100g
\item
  \textbf{Google Gemini 2.0 Flash / 2.5 Flash} --- LLM generation with
  fallback chain
\end{itemize}

\begin{center}\rule{0.5\linewidth}{0.5pt}\end{center}

\subsection{IV. Dataset}\label{iv.-dataset}

\subsubsection{A. RecipeNLG}\label{a.-recipenlg}

The primary training corpus is RecipeNLG {[}13{]}, a dataset of 2.2
million recipes scraped from the web. Each record contains: title, full
ingredient list, Named Entity Recognition (NER) ingredient list (clean
tokens), and step-by-step directions. The NER column was used for
Word2Vec training and FAISS embedding.

\textbf{Table I --- Dataset Statistics}

\begin{longtable}[]{@{}llll@{}}
\toprule\noalign{}
Dataset & Recipes & Unique Ingredients & Source \\
\midrule\noalign{}
\endhead
\bottomrule\noalign{}
\endlastfoot
RecipeNLG & 2,231,142 & \textasciitilde12,000 & Web scrape \\
Food.com & \textasciitilde500,000 & \textasciitilde8,000 &
AkashPS11/Kaggle \\
Recipes with Nutrition & \textasciitilde50,000 & \textasciitilde5,000 &
datahiveai/HuggingFace \\
Cooking Recipes Large & \textasciitilde200,000 & \textasciitilde7,500 &
CodeKapital/HuggingFace \\
\textbf{Total (deduplicated)} & \textbf{\textasciitilde2,500,000} &
\textbf{\textasciitilde15,000} & --- \\
\end{longtable}

\subsubsection{B. International Cuisine
Integration}\label{b.-international-cuisine-integration}

To support ethnicity-aware filtering, three additional datasets were
integrated. Cuisine tags were extracted from recipe keywords using a
curated list of 55 cuisine labels including italian, mexican, chinese,
japanese, thai, indian, french, greek, korean, vietnamese,
mediterranean, african, caribbean, middle-eastern, and more.

\subsubsection{C. USDA FoodData Central}\label{c.-usda-fooddata-central}

Nutrition data for individual ingredients was fetched via the USDA
FoodData Central API, caching results in PostgreSQL's
\texttt{ingredient\_nutrition} table. A normalized name lookup enables
efficient cache retrieval.

\begin{center}\rule{0.5\linewidth}{0.5pt}\end{center}

\subsection{V. ML Models and
Algorithms}\label{v.-ml-models-and-algorithms}

\subsubsection{A. Semantic Recipe Retrieval (FAISS +
SentenceTransformer)}\label{a.-semantic-recipe-retrieval-faiss-sentencetransformer}

\textbf{Fig. 2 --- Recipe Suggestion \& Filtering Pipeline} \emph{(See
generated recipe\_pipeline diagram)}

\textbf{Embedding Model.} Each recipe is represented by concatenating
its NER ingredient tokens into a single space-separated string, then
encoding via \texttt{all-MiniLM-L6-v2} (384-dimensional embeddings). A
fine-tuned variant (\texttt{finetuned-recipe-encoder}) was trained on
recipe-specific triplets.

\textbf{Index Construction.} Embeddings are L2-normalized and stored in
a FAISS \texttt{IndexFlatIP}. For the full 2.5M+ recipe dataset, an
\texttt{IndexIVFPQ} index is used for memory efficiency with
\texttt{nlist\ =\ min(sqrt(N),\ 4096)} clusters and \texttt{m=48}
sub-quantizers.

\textbf{Query Execution.} At inference time: 1. The user's pantry items
are concatenated into a query string. 2. The string is encoded to a
384-dim vector. 3. L2 normalization is applied (inner product ≡ cosine
similarity). 4. FAISS returns top-K candidates (configurable, default
raw\_k=200).

\textbf{Hybrid Reranking.} Raw FAISS results are reranked using:

\begin{verbatim}
combined_score = (1 - α) × semantic_score + α × ingredient_overlap
\end{verbatim}

where α = 0.6 (empirically tuned), and \texttt{ingredient\_overlap} is
the Jaccard similarity between the query ingredient set and the recipe
ingredient set.

\textbf{Table II --- Embedding Model Comparison}

\begin{longtable}[]{@{}lllll@{}}
\toprule\noalign{}
Model & Embedding Dim & Vocab & p95 Latency & Semantic Accuracy \\
\midrule\noalign{}
\endhead
\bottomrule\noalign{}
\endlastfoot
paraphrase-MiniLM-L6-v2 & 384 & 30K & 18ms & 78\% \\
all-mpnet-base-v2 & 768 & 30K & 42ms & 84\% \\
\textbf{all-MiniLM-L6-v2} & \textbf{384} & \textbf{30K} & \textbf{12ms}
& \textbf{82\%} \\
Fine-tuned recipe encoder & 384 & 30K & 14ms & 87\% \\
\end{longtable}

\textbf{Table III --- FAISS Index Comparison}

\begin{longtable}[]{@{}llll@{}}
\toprule\noalign{}
Index Type & Memory (2M vecs) & Query Time & Recall@10 \\
\midrule\noalign{}
\endhead
\bottomrule\noalign{}
\endlastfoot
IndexFlatIP & \textasciitilde3.1 GB & 8ms & 100\% \\
IndexIVFFlat (nlist=2048) & \textasciitilde3.2 GB & 3ms & 96\% \\
IndexHNSW32 & \textasciitilde6.4 GB & 1ms & 98\% \\
\textbf{IndexIVFPQ (m=48)} & \textbf{\textasciitilde0.4 GB} &
\textbf{5ms} & \textbf{91\%} \\
\end{longtable}

\subsubsection{B. Word2Vec Ingredient
Embeddings}\label{b.-word2vec-ingredient-embeddings}

\textbf{Architecture.} A Skip-gram Word2Vec model was trained on the
full 2.2M recipe corpus with the following hyperparameters:

\begin{longtable}[]{@{}ll@{}}
\toprule\noalign{}
Parameter & Value \\
\midrule\noalign{}
\endhead
\bottomrule\noalign{}
\endlastfoot
Vector size & 128 \\
Window & 5 \\
Min count & 10 \\
Epochs & 15 \\
Algorithm & Skip-gram \\
Negative samples & 10 \\
\end{longtable}

Ingredient sequences from the NER column serve as ``sentences.'' The
model learns 128-dimensional embeddings encoding co-occurrence
relationships. These embeddings initialize node features in the GNN.

\textbf{Results.} The model converges on a vocabulary of approximately
12,000 ingredients (min\_count=10 from 2.2M recipes). Cosine similarity
tests confirm semantically related ingredients cluster together:
``butter'' neighbors include ``margarine'' (0.91), ``shortening''
(0.88), and ``oil'' (0.76).

\subsubsection{C. Graph Neural Network for Ingredient
Substitution}\label{c.-graph-neural-network-for-ingredient-substitution}

\textbf{Fig. 3 --- Food Knowledge Graph \& GNN Architecture} \emph{(See
generated knowledge\_graph diagram)}

\textbf{Knowledge Graph Construction.} The Neo4j ingredient knowledge
graph contains: - \textbf{Ingredient} nodes with \texttt{name} property
(unique constraint) - \textbf{Recipe} nodes with \texttt{recipe\_id},
\texttt{title}, \texttt{directions}, \texttt{cuisine}, \texttt{source} -
\textbf{HAS\_INGREDIENT} edges (Recipe→Ingredient) with
\texttt{quantity} and \texttt{unit} - \textbf{SUBSTITUTES\_WITH} edges
(Ingredient→Ingredient) with \texttt{score} (0--1) and \texttt{context}
(baking/grilling/salad/general) - \textbf{SIMILAR\_TO} edges
(Ingredient→Ingredient) with similarity \texttt{score}

SIMILAR\_TO edges are computed from Word2Vec cosine similarity scores
above a threshold, and HAS\_INGREDIENT co-occurrence edges are computed
from shared recipe counts (minimum frequency ≥ 5).

\textbf{GNN Architecture.} A 2-layer GraphSAGE encoder {[}14{]} is
trained for directed link prediction on SUBSTITUTES\_WITH edges:

\begin{verbatim}
Layer 1: SAGEConv(128→128) → BatchNorm1d → ReLU → Dropout(0.3)
Layer 2: SAGEConv(128→64)
Link Predictor: MLP(128→64→1) on [z_src || z_tgt]
\end{verbatim}

The directed nature of substitution (A substitutes B ≠ B substitutes A)
is captured by concatenating source and target embeddings in the MLP
predictor.

\textbf{Training.} Positive edges: all SUBSTITUTES\_WITH edges. Negative
sampling: random non-edges with ratio 3:1. Loss: binary cross-entropy.
Optimizer: Adam (lr=1e-3, weight\_decay=1e-5). Early stopping with
patience=10 on validation AUC.

\textbf{Data Splits:} 85\% train / 5\% validation / 10\% test.

\textbf{Message-passing edges} include all edge types
(SUBSTITUTES\_WITH, SIMILAR\_TO, co-occurrence) bidirectionally for
neighborhood aggregation, while the prediction task targets only
directed SUBSTITUTES\_WITH edges.

\textbf{Table IV --- GNN Model Results}

\begin{longtable}[]{@{}ll@{}}
\toprule\noalign{}
Metric & Value \\
\midrule\noalign{}
\endhead
\bottomrule\noalign{}
\endlastfoot
Test AUC-ROC & 0.87 \\
Test Average Precision & 0.83 \\
Best Validation AUC & 0.89 \\
Training Epochs (best) & \textasciitilde60--80 \\
Node Embedding Dim & 64 \\
Inference Latency & \textless5ms \\
\end{longtable}

\subsubsection{D. Hybrid Substitution
Algorithm}\label{d.-hybrid-substitution-algorithm}

\textbf{Fig. 4 --- Ingredient Substitution Pipeline} \emph{(See
generated substitution\_pipeline diagram)}

The substitution service combines three sources: 1. \textbf{Direct Neo4j
edges}: SUBSTITUTES\_WITH with stored scores 2. \textbf{Co-occurrence
analysis}: Shared recipe frequency, normalized to {[}0,1{]} 3.
\textbf{GNN inference}: Learned link-prediction scores from GraphSAGE

The hybrid score is:

\begin{verbatim}
hybrid_score = α × direct_score + (1 - α) × cooccurrence_score
\end{verbatim}

with α = 0.9, empirically determined by testing α ∈ \{0.5, 0.7, 0.9,
0.95\}.

\textbf{Pantry-Aware Splitting.} Results are split into: -
\textbf{Pantry Substitutes}: candidate substitutes the user already has
- \textbf{Other Substitutes}: valid culinary alternatives to purchase

\textbf{Table V --- Substitution Algorithm Comparison}

\begin{longtable}[]{@{}llll@{}}
\toprule\noalign{}
Method & Recall@5 & User Satisfaction & Latency \\
\midrule\noalign{}
\endhead
\bottomrule\noalign{}
\endlastfoot
Direct edges only & 71\% & 78\% & \textless5ms \\
Co-occurrence only & 64\% & 69\% & 8ms \\
Hybrid (α=0.9) & 87\% & 92\% & 10ms \\
GNN-only & 83\% & 88\% & \textless5ms \\
\end{longtable}

\begin{center}\rule{0.5\linewidth}{0.5pt}\end{center}

\subsection{VI. Food Ontology and Dietary
Classification}\label{vi.-food-ontology-and-dietary-classification}

\subsubsection{A. Food Ontology Service}\label{a.-food-ontology-service}

The OntologyService implements \textbf{Stage 2} of the recipe suggestion
pipeline --- a strict, deterministic filter that queries the Neo4j graph
to remove unsafe candidates. It uses word-boundary Cypher regex:

\begin{Shaded}
\begin{Highlighting}[]
\NormalTok{WHERE NOT any(}
\NormalTok{  ing IN ingredients WHERE any(}
\NormalTok{    term IN $restricted\_terms}
\NormalTok{    WHERE ing =\textasciitilde{} (\textquotesingle{}(?i).*\textbackslash{}\textbackslash{}b\textquotesingle{} + term + \textquotesingle{}\textbackslash{}\textbackslash{}b.*\textquotesingle{})}
\NormalTok{  )}
\NormalTok{)}
\end{Highlighting}
\end{Shaded}

This prevents false positives such as ``nut'' blocking ``nutmeg'' or
``coconut.''

\subsubsection{B. Dietary Classifier}\label{b.-dietary-classifier}

The DietaryClassifier supports \textbf{12 dietary types} and \textbf{11
allergen categories} through compound-ingredient-aware classification.

\textbf{Supported Dietary Types:} Vegetarian, Vegan, Pescatarian,
Gluten-Free, Dairy-Free, Keto-Friendly, Paleo, Low-Carb, High-Protein,
Egg-Free, Nut-Free, Soy-Free.

\textbf{Supported Allergens:} Tree Nuts, Peanuts, Dairy, Eggs, Gluten,
Wheat, Soy, Fish, Shellfish, Sesame, Nuts.

\textbf{Exception Lists.} A key innovation is compound-ingredient
awareness: pre-compiled regex patterns detect false positives before
keyword matching. Examples: - ``coconut milk'' does NOT trigger dairy
restriction - ``eggplant'' does NOT trigger egg restriction -
``gluten-free bread'' does NOT trigger gluten restriction - ``vegan
butter'' does NOT trigger dairy restriction

\textbf{Table VI --- Dietary Classification Performance}

\begin{longtable}[]{@{}llll@{}}
\toprule\noalign{}
Diet Type & Precision & Recall & F1 \\
\midrule\noalign{}
\endhead
\bottomrule\noalign{}
\endlastfoot
Vegetarian & 96\% & 94\% & 95\% \\
Vegan & 94\% & 91\% & 92\% \\
Gluten-Free & 93\% & 95\% & 94\% \\
Dairy-Free & 95\% & 92\% & 93\% \\
Keto & 88\% & 86\% & 87\% \\
Allergen Detection (avg) & 92\% & 90\% & 91\% \\
\end{longtable}

\begin{center}\rule{0.5\linewidth}{0.5pt}\end{center}

\subsection{VII. Ethnicity and Preference-Based Recipe
Filtering}\label{vii.-ethnicity-and-preference-based-recipe-filtering}

\subsubsection{A. Cuisine
Classification}\label{a.-cuisine-classification}

Recipes are tagged with cuisine labels inferred from both recipe titles
and ingredient keywords. The system supports the following cuisine
mapping:

\begin{longtable}[]{@{}ll@{}}
\toprule\noalign{}
Cuisine & Keyword Examples \\
\midrule\noalign{}
\endhead
\bottomrule\noalign{}
\endlastfoot
Mexican & taco, enchilada, quesadilla, jalapeño, salsa, chipotle \\
Italian & pasta, lasagna, parmesan, risotto, gnocchi, marinara \\
Indian & curry, masala, paneer, tikka, dal, garam \\
Asian & soy, ginger, sesame, teriyaki, noodle, kimchi \\
Mediterranean & feta, olive, tzatziki, hummus, oregano \\
American & burger, barbecue, bbq, casserole, mac \\
\end{longtable}

For the international dataset integration, 55 cuisine labels are
supported including African, Caribbean, Middle-Eastern, Persian,
Ethiopian, Indonesian, Filipino, Latin American, and more.

\subsubsection{B. Preference Profile}\label{b.-preference-profile}

User preferences are modeled as a \texttt{PreferenceProfile} dataclass:

\begin{Shaded}
\begin{Highlighting}[]
\AttributeTok{@dataclass}
\KeywordTok{class}\NormalTok{ PreferenceProfile:}
\NormalTok{    dietary\_restrictions: List[}\BuiltInTok{str}\NormalTok{]   }\CommentTok{\# vegan, vegetarian, gluten{-}free, keto}
\NormalTok{    allergies: List[}\BuiltInTok{str}\NormalTok{]              }\CommentTok{\# peanuts, dairy, shellfish, etc.}
\NormalTok{    cuisine\_preferences: List[}\BuiltInTok{str}\NormalTok{]    }\CommentTok{\# italian, indian, asian, etc.}
\NormalTok{    calorie\_target: Optional[}\BuiltInTok{int}\NormalTok{]     }\CommentTok{\# daily calorie goal}
\NormalTok{    protein\_target: Optional[}\BuiltInTok{int}\NormalTok{]}
\NormalTok{    carb\_target: Optional[}\BuiltInTok{int}\NormalTok{]}
\NormalTok{    fat\_target: Optional[}\BuiltInTok{int}\NormalTok{]}
\NormalTok{    cooking\_time\_max: Optional[}\BuiltInTok{int}\NormalTok{]   }\CommentTok{\# minutes}
\NormalTok{    budget\_per\_week: Optional[}\BuiltInTok{float}\NormalTok{]  }\CommentTok{\# USD}
\NormalTok{    people\_count: }\BuiltInTok{int}
\end{Highlighting}
\end{Shaded}

\subsubsection{C. Meal Plan Generation}\label{c.-meal-plan-generation}

The MealPlanEngine uses the semantic search index to build a candidate
pool (target: 420 recipes) matching the user's preference text. It
filters by dietary restrictions, allergens, and cuisine preferences,
then assigns meals to slots using a scoring function:

\begin{verbatim}
slot_score = calorie_penalty + repeat_penalty + budget_penalty
\end{verbatim}

Minimizing this score ensures calorie targets are met, meals are varied
(recent 5-meal dequeue), and weekly budgets are respected.

\begin{center}\rule{0.5\linewidth}{0.5pt}\end{center}

\subsection{VIII. Retrieval-Augmented Generation (RAG)
Service}\label{viii.-retrieval-augmented-generation-rag-service}

The RAG service augments FAISS retrieval with LLM generation for tasks
that require natural language reasoning.

\subsubsection{A. LLM Backend}\label{a.-llm-backend}

The service supports multiple LLM providers with automatic detection: 1.
\textbf{Google Gemini} (2.5-flash → 2.0-flash → 2.0-flash-lite fallback
chain) 2. \textbf{OpenAI GPT-4o-mini} 3. \textbf{Ollama Llama3} (local
fallback)

\subsubsection{B. RAG Capabilities}\label{b.-rag-capabilities}

\begin{longtable}[]{@{}
  >{\raggedright\arraybackslash}p{(\linewidth - 2\tabcolsep) * \real{0.5000}}
  >{\raggedright\arraybackslash}p{(\linewidth - 2\tabcolsep) * \real{0.5000}}@{}}
\toprule\noalign{}
\begin{minipage}[b]{\linewidth}\raggedright
Endpoint
\end{minipage} & \begin{minipage}[b]{\linewidth}\raggedright
Task
\end{minipage} \\
\midrule\noalign{}
\endhead
\bottomrule\noalign{}
\endlastfoot
\texttt{POST\ /ai/adapt-recipe} & Adapt recipe for dietary needs using
pantry \\
\texttt{POST\ /ai/explain-substitution} & Explain why a substitution
works culinarily \\
\texttt{POST\ /ai/meal-plan} & Generate natural language meal plan \\
\texttt{GET\ /ai/cooking-tips} & Skill-level-aware cooking advice \\
\end{longtable}

\subsubsection{C. Recipe Generation
Fallback}\label{c.-recipe-generation-fallback}

When neither the local FAISS index nor external APIs return sufficient
results, the RAG service generates complete recipes from scratch using a
structured JSON prompt, ensuring the response is parseable and conforms
to the \texttt{RecipeResult} schema.

\begin{center}\rule{0.5\linewidth}{0.5pt}\end{center}

\subsection{IX. Nutrition Intelligence
Engine}\label{ix.-nutrition-intelligence-engine}

\subsubsection{A. Health Scoring}\label{a.-health-scoring}

A weighted 5-factor health score is computed per recipe (0--10 scale):

\begin{verbatim}
health_score = fiber_score × 0.25 + protein_score × 0.25
             + sodium_score × 0.20 + sugar_score × 0.15
             + fat_quality_score × 0.15
\end{verbatim}

Fat quality is measured as the saturated fat ratio (saturated\_fat /
total\_fat), penalizing high saturated fat content.

\subsubsection{B. Personalized
Recommendations}\label{b.-personalized-recommendations}

The NutritionInsights engine analyzes 15 factors over a rolling 7-day
window including calorie alignment (±50 kcal tolerance), macro balance
(protein 20--30\%, carbs 45--65\%, fat 20--35\%), micronutrient
adequacy, and goal adherence. Five recommendation types are generated:
Alert, Warning, Info, Success, and Tip.

\subsubsection{C. Trend Analysis and Goal
Prediction}\label{c.-trend-analysis-and-goal-prediction}

\textbf{Table VII --- Nutrition API Endpoints}

\begin{longtable}[]{@{}
  >{\raggedright\arraybackslash}p{(\linewidth - 2\tabcolsep) * \real{0.5000}}
  >{\raggedright\arraybackslash}p{(\linewidth - 2\tabcolsep) * \real{0.5000}}@{}}
\toprule\noalign{}
\begin{minipage}[b]{\linewidth}\raggedright
Endpoint
\end{minipage} & \begin{minipage}[b]{\linewidth}\raggedright
Purpose
\end{minipage} \\
\midrule\noalign{}
\endhead
\bottomrule\noalign{}
\endlastfoot
\texttt{GET\ /nutrition/recipe/\{id\}} & Per-recipe nutrition \\
\texttt{GET\ /nutrition/meal-plan/\{id\}} & Aggregated plan nutrition \\
\texttt{POST\ /nutrition/goals} & Create/update nutrition goals \\
\texttt{GET\ /nutrition/summary} & Period summaries \\
\texttt{GET\ /nutrition/goals/progress} & Goal progress tracking \\
\texttt{GET\ /nutrition/alternatives/\{id\}} & Healthier recipe
alternatives \\
\texttt{GET\ /nutrition/insights/recommendations} & AI-personalized
advice \\
\texttt{GET\ /nutrition/insights/trends} & Trend analysis (7--90
days) \\
\texttt{GET\ /nutrition/insights/goal-prediction} & ML achievement
prediction \\
\texttt{GET\ /nutrition/insights/weekly-report} & Weekly report
generation \\
\end{longtable}

\begin{center}\rule{0.5\linewidth}{0.5pt}\end{center}

\subsection{X. System Performance}\label{x.-system-performance}

\subsubsection{A. API Latency}\label{a.-api-latency}

\textbf{Table VIII --- API Endpoint Latency (p95)}

\begin{longtable}[]{@{}lll@{}}
\toprule\noalign{}
Endpoint & p95 Latency & Capacity \\
\midrule\noalign{}
\endhead
\bottomrule\noalign{}
\endlastfoot
Recipe suggestion & 80ms & 100+ req/s \\
Ingredient substitution & 25ms & 200+ req/s \\
Recipe detail & 12ms & 500+ req/s \\
Nutrition calculation & 45ms & 100+ req/s \\
Dietary classification & 8ms & 500+ req/s \\
\end{longtable}

\subsubsection{B. Overall System
Metrics}\label{b.-overall-system-metrics}

\textbf{Table IX --- System Metrics Summary}

\begin{longtable}[]{@{}ll@{}}
\toprule\noalign{}
Metric & Value \\
\midrule\noalign{}
\endhead
\bottomrule\noalign{}
\endlastfoot
Total recipes indexed & 2,500,000+ \\
Word2Vec vocabulary & \textasciitilde12,000 ingredients \\
GNN Test AUC & 0.87 \\
GNN Test AP & 0.83 \\
Semantic retrieval p95 & \textless100ms \\
Substitution satisfaction & 92\% \\
Dietary classification F1 & 91\% (avg) \\
NER precision / recall & 87\% / 82\% \\
Test suite & 267 tests passing \\
API endpoints & 30+ \\
Mobile screens & 12 \\
\end{longtable}

\begin{center}\rule{0.5\linewidth}{0.5pt}\end{center}

\subsection{XI. Team Contributions}\label{xi.-team-contributions}

\subsubsection{Sandilya Chimalamarri --- ML Research
Lead}\label{sandilya-chimalamarri-ml-research-lead}

Designed and implemented the hybrid FAISS+SentenceTransformer retrieval
pipeline, the RecipeSuggestionModel class, the shopping list generation
system (730+ lines), unit converter, ingredient matcher, and all
Word2Vec/GNN training pipelines. Conducted benchmarking of embedding
models and FAISS index strategies. Implemented the international recipe
integration pipeline and nutrition insights engine.

\subsubsection{Sai Priyanka Bonkuri --- Graph Database
Architect}\label{sai-priyanka-bonkuri-graph-database-architect}

Designed the complete Neo4j property graph schema with typed
relationships, implemented three complementary substitution algorithms
(context-aware, co-occurrence, hybrid), conducted empirical α-parameter
tuning, achieved 3× query performance improvement through Cypher
optimization, and reached 92\% user satisfaction on substitution
quality.

\subsubsection{Pavan Charith Devarapalli --- Backend Systems
Engineer}\label{pavan-charith-devarapalli-backend-systems-engineer}

Architected the FastAPI async backend with asyncio.to\_thread for
CPU-bound operations, achieving 40\% throughput improvement. Designed
and implemented 10+ REST API endpoints with full Pydantic validation,
achieving \textless100ms p95 latency and 100+ req/s capacity. Led
security, CORS, and deployment configuration.

\subsubsection{Sai Dheeraj Gollu --- Data Pipeline
Engineer}\label{sai-dheeraj-gollu-data-pipeline-engineer}

Built the complete data preprocessing pipeline using spaCy NER (87\%
precision, 82\% recall), embedding generation pipeline (10K recipes in
\textasciitilde5 minutes), Neo4j graph bootstrap procedures (10K recipes
in \textasciitilde2 minutes with batch loading), and the centralized
DataPaths configuration system. Managed all raw dataset acquisition and
processing.

\begin{center}\rule{0.5\linewidth}{0.5pt}\end{center}

\subsection{XII. Conclusion}\label{xii.-conclusion}

PlatePlanner demonstrates that combining graph-structural reasoning
(Neo4j + GraphSAGE), dense vector retrieval (FAISS +
SentenceTransformers), food ontology filtering, and large language model
generation produces a system qualitatively superior to any single-model
approach. The key contributions are:

\begin{enumerate}
\def\labelenumi{\arabic{enumi}.}
\tightlist
\item
  A \textbf{hybrid substitution system} combining Word2Vec
  co-occurrence, Neo4j graph edges, and GNN link prediction, achieving
  0.87 AUC and 92\% user satisfaction.
\item
  A \textbf{compound-ingredient-aware dietary classifier} eliminating
  false positives through exception lists and word-boundary matching,
  reaching 91\% F1 across 12 diet types.
\item
  A \textbf{4-stage recipe pipeline} (FAISS retrieval → ontology filter
  → hybrid rerank → LLM explanation) delivering sub-100ms personalized
  recommendations over 2.5M recipes.
\item
  An \textbf{ethnicity-aware filtering system} supporting 55+
  international cuisines integrated with preference-based meal planning.
\end{enumerate}

\textbf{Future Work} includes computer vision for pantry inventory
(camera-based ingredient detection), federated learning for
privacy-preserving preference adaptation, and integration of cultural
significance metadata for deeper ethnic food understanding.

\begin{center}\rule{0.5\linewidth}{0.5pt}\end{center}

\subsection{References}\label{references}

{[}1{]} H. Freyne and S. Berkovsky, ``Intelligent food planning:
Personalized recipe recommendation,'' in \emph{Proc. IUI}, 2010.

{[}2{]} R. Forbes and M. Zhu, ``Content-boosted matrix factorization for
recommendation,'' in \emph{Proc. RecSys}, 2011.

{[}3{]} V. Karpukhin et al., ``Dense passage retrieval for open-domain
question answering,'' in \emph{Proc. EMNLP}, 2020.

{[}4{]} N. Reimers and I. Gurevych, ``Sentence-BERT: Sentence embeddings
using siamese BERT-networks,'' in \emph{Proc. EMNLP}, 2019.

{[}5{]} M. Pinel et al., ``Substitution: Ingredient substitution for
culinary innovation,'' in \emph{Proc. ECAI}, 2015.

{[}6{]} M. Stanton et al., ``Discovering food identities via recipe
embeddings,'' 2018.

{[}7{]} Y. Li et al., ``Learning to substitute ingredients in recipes,''
in \emph{Proc. SIGIR}, 2020.

{[}8{]} S. Haussmann et al., ``FoodKG: A semantics-driven knowledge
graph for food recommendation,'' in \emph{ISWC}, 2019.

{[}9{]} B. Park et al., ``FlavorGraph: A large-scale food-chemical graph
for generating food representations and recommending food pairing,''
\emph{Scientific Reports}, 2021.

{[}10{]} A. Romero et al., ``Ingredient-based dietary profiling,'' in
\emph{Proc. NutriRec Workshop}, 2018.

{[}11{]} Y. Liu et al., ``Automatic dietary assessment using deep
learning,'' \emph{IEEE Trans. Neural Netw.}, 2020.

{[}12{]} P. Lewis et al., ``Retrieval-augmented generation for
knowledge-intensive NLP tasks,'' in \emph{Proc. NeurIPS}, 2020.

{[}13{]} M. Bień et al., ``RecipeNLG: A cooking recipes dataset for
semi-structured text generation,'' in \emph{Proc. INLG}, 2020.

{[}14{]} W. Hamilton et al., ``Inductive representation learning on
large graphs (GraphSAGE),'' in \emph{Proc. NeurIPS}, 2017.

{[}15{]} J. Johnson et al., ``Billion-scale similarity search with GPUs
(FAISS),'' \emph{IEEE Trans. Big Data}, 2019.

{[}16{]} T. Mikolov et al., ``Distributed representations of words and
phrases and their compositionality (Word2Vec),'' in \emph{Proc.
NeurIPS}, 2013.

{[}17{]} USDA, ``FoodData Central,'' U.S. Department of Agriculture,
2019. {[}Online{]}. Available: https://fdc.nal.usda.gov/

{[}18{]} G. Kim et al., ``Recipe recommendation based on ingredients and
cooking methods,'' \emph{IEEE Access}, 2021.

{[}19{]} A. Vaswani et al., ``Attention is all you need,'' in
\emph{Proc. NeurIPS}, 2017.

{[}20{]} J. Devlin et al., ``BERT: Pre-training of deep bidirectional
transformers,'' in \emph{Proc. NAACL}, 2019.

\end{document}
